
\documentclass{book}

\usepackage{
  draftwatermark, % diagonal DRAFT watermark
  enumerate,      % enumerate with letter labels
  fancyhdr,       % custom chapter headings
  fullpage,       % reasonable margins
  hyperref,       % hotlinks
  makeidx,        % index
  tcolorbox,      % background colors in sidebars
  times,          % font choice
}

% Don't say "Chapter 3", just "3"
\renewcommand{\chaptername}{}

% do not indent paragraphs horizontally
\setlength{\parindent}{0pt}

% blank line between paragraphs
\setlength{\parskip}{12pt}

% draft watermark
\SetWatermarkText{DRAFT}
\SetWatermarkColor[gray]{0.9}

% name of our campus, so this is consistent everywhere
\newcommand{\campusname}{CSUF}

% tip sidebar
\newenvironment{tip}
    {\tcolorbox \begin{tabular}{p{2cm} l} \Large{TIP} & }
    {\end{tabular} \endtcolorbox}

\makeindex

\begin{document}

\title{Undergraduate Handbook \\ 2015-2016 Edition}
\author{Department of Computer Science \\ California State University, Fullerton}
\date{Draft: \today}
\maketitle

\newpage
\tableofcontents

\chapter{Introduction}

\section{The field of computer science}
Computer Science is the systematic study of computing systems and computation. The body of knowledge contains the theoretical foundation for understanding computing systems and methods, design methodology, algorithms, and software and hardware tools.

These programs cover a wide range of areas, including:
\begin{itemize}
\item multimedia and digital game technologies,
\item Internet and enterprise computing,
\item wireless and mobile computing,
\item databases and data mining,
\item computer security,
\item software engineering, and
\item computational bioinformatics.
\end{itemize}

Computer Science prepares graduates for rewarding careers in all areas of business, government, education and industry. These organizations, large and small, need computer professionals to address their needs with specific programs and systems. Computer science professionals tackle complicated problems and create computer solutions to solve them, devising new ways to use computers.

\section{The department}
\label{section:the_department}
The faculty and staff of the Computer Science Department welcome you into our program and sincerely wish you good luck on your journey into higher education, and continued success.

Whenever you have a question about the Department---its policies, its curriculum, its services, your progress, or anything else---feel free to contact us.

\index{contact information}
\begin{tabular}{rp{6in}} % right column in p mode so postal address can have newlines
  Web: & \href{http://fullerton.edu/ecs/cs/}{http://fullerton.edu/ecs/cs/} \\ \index{website}
  E-mail: & \href{mailto:csoffice@ecs.fullerton.edu}{csoffice@ecs.fullerton.edu} \\ \index{e-mail}
  In person: & Room CS-522 \\ \index{department office}
  Telephone: & (657) 278-3700 \\ \index{phone number} \index{telephone number}
  Fax: & (657) 278-7168 \\ \index{fax number}
  Postal mail: & California State University, Fullerton \newline \index{postal address} \index{address}
Department of Computer Science \newline
P.O. Box 6870 \newline
Fullerton, CA 92834-6870
\end{tabular}

\section{Accreditation} \index{accreditation}\index{ABET}

The Bachelor of Science in Computer Science degree at \campusname is accredited by the Computer Accreditation Commission of ABET (\href{http://www.abet.org}{abet.org}).

\begin{center}
  \includegraphics{CAC-CMYK-W.pdf}
\end{center}

\section{The programs}

The Department offers the following Undergraduate programs, which are documented in this Handbook:
\begin{enumerate}
\item Bachelor of Science in Computer Science (BS CS), and
\item Minor in Computer Science.
\end{enumerate}

The Department also offers Graduate programs, which are documented elsewhere:
\begin{enumerate}
\item Master of Science in Computer Science (MS CS),
\item Master of Science in Software Engineering (MSE), and
\item Accelerated Master of Science in Software Engineering (AMSE).
\end{enumerate}

CS courses are also components of Computer Engineering, Electrical Engineering, and Mathematics programs at \campusname.

\section{Objectives and Outcomes}
\index{Program Educational Objectives (PEOs)} \index{Program Outcomes}
The Program Educational Objectives and Program Outcomes for the BS CS are \href{http://catalog.fullerton.edu/preview_entity.php?catoid=2&ent_oid=113&returnto=137#Learning_Goals_and_Student_Learning_Outcomes}{documented in the University Catalog}.

\section{Using this document}

This handbook covers information on how to earn a Bachelor of Science or a Minor in Computer Science, and contains information relevant to students pursuing them. If you are pursuing a Masters degree, please trade this handbook for a copy of the Graduate Handbook.

Some aspects of our programs are complex, and you may find it difficult to choose among alternatives. In those cases, we present our suggested default choice as a tip, as shown below. You are not required to obey these tips, but doing so is often a wise choice.
\begin{tip} \index{tips}
When in doubt, heed tips such as this one.
\end{tip}

\chapter{Other sources of information}

\chapter{The CS Major}

\section{Major requirements at a glance}

\section{Major prerequisite tree}

\section{Lower Division Core}
\section{Mathematics Core}
\section{Science and Mathematics Electives}
\section{Examination in Programming Proficiency (EPP)}
\section{Upper Division Core}
\section{Elective Tracks}
\section{Free Electives}
\label{section:free_electives}

Your Elective Track will require you to take some number of Free Elective courses. You may need to take additional Free Electives if you are short on units due to the Placement Examination, transfer, or other circumstances.

A course may be used as a Free Elective if it is
\begin{enumerate}
\item a CPSC course;
\item 3 units;
\item upper-division (300/400-level); \emph{and}
\item a requirement for any of the Elective Tracks.
\end{enumerate}

Therefore, the following courses may count as Free Electives:
\begin{itemize}
\item CPSC 303 - Multimedia Concepts (3)
\item CPSC 322L - Introduction to Computer-Aided Design (3)
\item CPSC 353 - Introduction to Computer Security (3)
\item CPSC 376 - Client/Server Systems with Java (3)
\item CPSC 386 - Introduction to Game Design and Production (3)
\item CPSC 439 - Theory of Computation (3)
\item CPSC 452 - Cryptography (3)
\item CPSC 454 - Cloud Computing and Security (3)
\item CPSC 456 - Network Security Fundamentals (3)
\item CPSC 459 - Micro-Computer Software Systems (3)
\item CPSC 462 - Software Design (3)
\item CPSC 463 - Software Testing (3)
\item CPSC 464 - Software Architecture (3)
\item CPSC 466 - Software Process (3)
\item CPSC 477 - Introduction to Grid Computing (3)
\item CPSC 483 - Data Mining and Pattern Recognition (3)
\item CPSC 484 - Principles of Computer Graphics (3)
\item CPSC 485 - Computational Bioinformatics (3)
\item CPSC 486 - Game Programming (3)
\item CPSC 489 - Game Development Project (3)
\item CPSC 491T - Variable Topics in Computer Science (3)
\item CPSC 495 - Internship in Computer Science (3)
\item CPSC 499 - Independent Study (3)
\end{itemize}

You may be able to use an adviser-approved course not on this list as a Free Elective. Such a course must be at least 3 units and directly related to your academic goals. If this interests you, discuss it with a major adviser. You may need to file a petition with the Department, which may or may not be approved.

\section{Upper division writing requirement}
\section{General Education (GE)}
\section{Academic Requirements}
\section{Double-majoring in a related field}

\chapter{The CS Minor}

\section{Minor requirements at a glance}
\section{Minor prerequisite tree}
\section{Suggested minor electives}
\section{For majors in related fields}

\chapter{Major tracks}

\section{Multimedia and Digital Games (MG)}
\section{Internet and Enterprise Computing (IE)}
\section{Software Engineering (SE)}
\section{Scientific Computing (SC)}
\section{Customized (CT)}

\chapter{Alternative pathways}

\section{Transfer}

\section{Computer Science Placement Examination}
\index{placement examination}
Our courses CPSC 120 and CPSC 121 cover computer programming. If you are proficient in this material, but do not have academic credit, you may attempt to establish your proficiency and skip one or both courses by taking the Computer Science Placement Examination. If you have taken CPSC 120, 121, or 131 at \campusname, you are not eligible to take the Placement Examination. See the \href{http://catalog.fullerton.edu/preview_program.php?catoid=2&poid=537&returnto=137}{Department's section of the University Catalog} for more information.

If you test out of one or both courses, you will be short on units in the Lower Division Core major requirement. You will need to earn substitute units by taking extra units of CPSC 223 and/or Upper Divsion Elective Courses (listed in section \ref{section:free_electives}.

\section{Advanced Placement (AP)}
\index{advanced placement (AP)}
If you took the Computer Science AP exam and scored well, you may be able to get credit for CPSC 120, or both 120 and 121. See the \href{http://catalog.fullerton.edu/content.php?catoid=2&navoid=153#Credit%20by%20Advanced%20Placement%20Chart}{Credit by Advanced Placement Chart} in the University Catalog.

\section{Internships}

\section{Independent Study}

\chapter{Advisement}

It can be frustrating to find out that you took a class that wasn’t useful for your course of study. Not being able to take a class when you want because of a needed prerequisite is even worse---it slows your progress and can delay your graduation. To avoid problems like these, the University offers advisement counseling to all students. This is your opportunity to review your progress toward your degree and to discuss electives that match your career goals.

\section{Major Advisement}
\index{advisement!major}
You have to set up an advisement appointment yourself. Contact the department and ask for an advisement appointment. Our contact information is in section \ref{section:the_department}

\section{Required Advisement}
\index{advisement!required}
The College of Engineering and Computer Science places a registration hold on all undergraduate students once a year to ensure the student meets with a department advisor. Students whose surname begins with ``A'' thru ``M'' will have a hold each Fall term. Those whose surnname begins with ``N'' thru ``Z'' will find a registration hold each Spring term.  You will not be able to register for any courses until you consult with a department advisor and the hold is subsequently removed.

\section{General Education (GE) Advisement}
\index{advisement!GE} \index{general education}
The University encourages all students to seek GE advisement, each semester, well in advance of registration.  You may obtain information about the Cal State Fullerton GE curriculum and degree requirements by visiting the Academic Advisement Center in UH-123B.

\section{First-time freshmen}
\index{freshmen}
You should make an appointment to see the department adviser as early as possible. It’s very important that you understand the program and the sequence in which you should take courses.

\section{Transfer students}
\index{transfer}
You should make an advisement appointment as early as possible. The department adviser can answer your questions about transfer credit for general education courses and can evaluate courses that apply to your major. Please bring any transcripts or grade reports you have, official or not, to this appointment. A catalog from your prior institution may prove useful, particularly from those outside the Orange County area.

\section{Nearing graduation (within one year)}
\index{graduation}
After completing 90 units of coursework, you are eligible to apply for graduation. The only way to apply for graduation is online through the TITAN Online Student Center. You cannot graduate without a completed Grad Check. \index{grad check} The University Catalog \href{http://catalog.fullerton.edu/content.php?catoid=2&navoid=108#Graduation_Requirement_Check}{has more information on Grad Checks}.

\section{Probation}
\index{probation}
If you are on probation, it is definitely time to see an adviser. Until you do so, a hold will be in place on your file, preventing you from registering in classes. Your adviser will discuss with you the problems that led to your probation and review strategies you should take to get off probation. Make your advisement appointment early so your registration is not held up.

\section{Your Catalog Year}
\index{catalog year}
Be sure to follow the course requirements for your catalog year. Your catalog year is determined by the Admissions Office and is a part of your student records. Typically, this is the year you began college; occasionally an adviser may approve a later year.

\chapter{Resources and Activities}

\section{Open Labs}

\section{Tutoring Center}

\section{Supplemental Instruction}

\section{Clubs}
\subsection{Association of Computing Machinery (ACM)}
\index{ACM}

\subsection{ACM-W}
\index{ACM}\index{ACM-W}

\subsection{Upsilon Pi Epsilon (UPE)}
\subsection{Video Game Design Club (VGDC)}
\subsection{Security Club (?)}

\chapter{Progress Flowcharts}

\chapter{Credits and Revision History}

Copyright 2015, Department of Computer Science, California State University, Fullerton.

% magic to add index to table of contents without being a numbered chapter, from
% https://en.wikibooks.org/wiki/LaTeX/Indexing#Adding_index_to_table_of_contents
\cleardoublepage \addcontentsline{toc}{chapter}{Index} \printindex

\end{document}
